\section{Telemetría: el servicio notario}

El notario es un servicio Python independiente que monitoriza en tiempo real el estado
de la sesión de Voctomix. Se suscribe al protocolo de control TCP de voctocore como
un cliente más y registra todos los cambios de estado: fuente de vídeo activa, modo
de composite, estado del stream blanker, overlays cargados, etc.

\subsection{Exportación de datos}

Los datos se exportan en dos formatos complementarios:

\begin{itemize}
  \item \textbf{API REST HTTP (puerto 8080):} devuelve el estado actual de la sesión
        en formato JSON, accesible por cualquier cliente HTTP. Resulta útil para
        dashboards de monitorización en tiempo real.
  \item \textbf{Mensajería AMQP (RabbitMQ):} cada cambio de estado se publica como
        un mensaje en el broker RabbitMQ, permitiendo la integración con sistemas
        de registro, alertas o automatización mediante colas de mensajes.
\end{itemize}

\subsection{Diseño no-SPOF}

La arquitectura desacoplada del notario ---conectado como un cliente TCP ordinario,
sin privilegios especiales sobre el núcleo--- garantiza que su eventual fallo no
afecte al procesamiento multimedia principal (\textit{single point of failure} libre).
El sistema de producción puede operar con plenas garantías aunque el servicio de
telemetría no esté disponible.

El módulo \texttt{voctogui/lib/gui\_state\_exporter.py} complementa esta funcionalidad
desde el lado de la GUI, serializando el estado de la interfaz en ficheros JSON que
permiten retomar una sesión interrumpida sin pérdida de configuración.

% --- Figura sugerida ---
% \begin{figure}[H]
%   \centering
%   \includegraphics[width=0.85\textwidth]{figuras/integracion_notario.png}
%   \caption{Diagrama de integración: voctocore -- notario -- RabbitMQ -- consumidores.}
%   \label{fig:integracion_notario}
% \end{figure}

% --- Tabla sugerida ---
% \begin{table}[H]
%   \centering
%   \caption{Endpoints REST del servicio notario.}
%   \label{tab:endpoints_notario}
%   \begin{tabular}{|l|c|l|}
%     \hline
%     \textbf{Ruta} & \textbf{Método} & \textbf{Descripción} \\
%     \hline
%     \texttt{/status}      & GET & Estado completo de la sesión en JSON \\
%     \texttt{/video}       & GET & Fuente de vídeo activa (A y B) \\
%     \texttt{/composite}   & GET & Modo de composite activo \\
%     \texttt{/stream}      & GET & Estado del stream blanker \\
%     \hline
%   \end{tabular}
% \end{table}
